\documentclass[11pt]{article}
\usepackage[margin=2.2cm]{geometry}
\usepackage{listings,xcolor}
\usepackage[most]{tcolorbox}
\usepackage{lmodern}
\lstset{basicstyle=\ttfamily\small,keywordstyle=\color{blue!70!black}\bfseries,commentstyle=\color{green!50!black},backgroundcolor=\color{gray!10},frame=single,breaklines=true}
\newtcolorbox{tipbox}[1]{colback=yellow!10,colframe=yellow!60!black,coltitle=black,fonttitle=\bfseries,title={#1}}
\newtcolorbox{notebox}{colback=cyan!5,colframe=cyan!50!black,enhanced,borderline west={4pt}{0pt}{cyan!50!black}}
\begin{document}
\section{Stacks}
A stack supports two operations: \texttt{push(x)} adds an element on top, and \texttt{pop()} removes and returns the top element. The function \texttt{is\_empty(S)} returns true if the stack $S$ has no elements.
\begin{lstlisting}[language=Python]
class Stack:
    def __init__(self):
        self.items = []          # the underlying list

    def push(self, x):
        self.items.append(x)

    def pop(self):
        if self.is_empty():
            raise IndexError("pop from empty stack")
        return self.items.pop()

    def is_empty(self):
        return len(self.items) == 0
\end{lstlisting}
\begin{tipbox}{Tip}
Use \texttt{collections.deque} when you need fast appends and pops from both ends. Calling \texttt{list.pop(0)} takes $O(n)$ time.
\end{tipbox}
\section{Linked lists in C}
Each node stores a value and a pointer to the next node:
\begin{lstlisting}[language=C]
struct node {
    int value;
    struct node *next;
};

int length(struct node *head) {
    int n = 0;
    for (struct node *p = head; p != NULL; p = p->next)
        n++;
    return n;
}
\end{lstlisting}
\begin{notebox}
\textbf{Note.} If \texttt{head == NULL} the list is empty and \texttt{length} returns 0. Always check for \texttt{NULL} before dereferencing a pointer.
\end{notebox}
\section{Algorithm: binary search}
\begin{tcolorbox}[colback=white,colframe=red!60!black,title={Algorithm 1: BinarySearch(A, x)}]
\begin{tabbing}
\quad\=\quad\=\quad\=\kill
1: $lo \gets 0$, $hi \gets n - 1$\\
2: \textbf{while} $lo \le hi$ \textbf{do}\\
3: \> $mid \gets \lfloor (lo + hi)/2 \rfloor$\\
4: \> \textbf{if} $A[mid] = x$ \textbf{then return} $mid$\\
5: \> \textbf{else if} $A[mid] < x$ \textbf{then} $lo \gets mid + 1$\\
6: \> \textbf{else} $hi \gets mid - 1$\\
7: \textbf{return} $-1$
\end{tabbing}
\end{tcolorbox}
Binary search runs in $O(\log n)$ time because the search interval halves at each step. The loop invariant is: if $x$ is in $A$, then $x \in A[lo..hi]$.
\begin{center}
\begin{tabular}{lll}
\textbf{Command} & \textbf{Meaning} & \textbf{Example}\\
\texttt{git status} & show changed files & \texttt{git status -s}\\
\texttt{git add} & stage changes & \texttt{git add main.py}\\
\texttt{git commit} & record a snapshot & \texttt{git commit -m "fix"}\\
\texttt{ls -la} & list all files & \texttt{ls -la /tmp}\\
\end{tabular}
\end{center}
\section{Exercises}
\begin{enumerate}
\item Write a function \texttt{reverse(head)} that reverses a linked list in place.
\item What is the output of \texttt{print([x*x for x in range(5)])}?
\item Explain why \texttt{a == b} and \texttt{a is b} can give different results in Python.
\end{enumerate}


\section{Stacks}
A stack supports two operations: \texttt{push(x)} adds an element on top, and \texttt{pop()} removes and returns the top element. The function \texttt{is\_empty(S)} returns true if the stack $S$ has no elements.
\begin{lstlisting}[language=Python]
class Stack:
    def __init__(self):
        self.items = []          # the underlying list

    def push(self, x):
        self.items.append(x)

    def pop(self):
        if self.is_empty():
            raise IndexError("pop from empty stack")
        return self.items.pop()

    def is_empty(self):
        return len(self.items) == 0
\end{lstlisting}
\begin{tipbox}{Tip}
Use \texttt{collections.deque} when you need fast appends and pops from both ends. Calling \texttt{list.pop(0)} takes $O(n)$ time.
\end{tipbox}
\section{Linked lists in C}
Each node stores a value and a pointer to the next node:
\begin{lstlisting}[language=C]
struct node {
    int value;
    struct node *next;
};

int length(struct node *head) {
    int n = 0;
    for (struct node *p = head; p != NULL; p = p->next)
        n++;
    return n;
}
\end{lstlisting}
\begin{notebox}
\textbf{Note.} If \texttt{head == NULL} the list is empty and \texttt{length} returns 0. Always check for \texttt{NULL} before dereferencing a pointer.
\end{notebox}
\section{Algorithm: binary search}
\begin{tcolorbox}[colback=white,colframe=red!60!black,title={Algorithm 1: BinarySearch(A, x)}]
\begin{tabbing}
\quad\=\quad\=\quad\=\kill
1: $lo \gets 0$, $hi \gets n - 1$\\
2: \textbf{while} $lo \le hi$ \textbf{do}\\
3: \> $mid \gets \lfloor (lo + hi)/2 \rfloor$\\
4: \> \textbf{if} $A[mid] = x$ \textbf{then return} $mid$\\
5: \> \textbf{else if} $A[mid] < x$ \textbf{then} $lo \gets mid + 1$\\
6: \> \textbf{else} $hi \gets mid - 1$\\
7: \textbf{return} $-1$
\end{tabbing}
\end{tcolorbox}
Binary search runs in $O(\log n)$ time because the search interval halves at each step. The loop invariant is: if $x$ is in $A$, then $x \in A[lo..hi]$.
\begin{center}
\begin{tabular}{lll}
\textbf{Command} & \textbf{Meaning} & \textbf{Example}\\
\texttt{git status} & show changed files & \texttt{git status -s}\\
\texttt{git add} & stage changes & \texttt{git add main.py}\\
\texttt{git commit} & record a snapshot & \texttt{git commit -m "fix"}\\
\texttt{ls -la} & list all files & \texttt{ls -la /tmp}\\
\end{tabular}
\end{center}
\section{Exercises}
\begin{enumerate}
\item Write a function \texttt{reverse(head)} that reverses a linked list in place.
\item What is the output of \texttt{print([x*x for x in range(5)])}?
\item Explain why \texttt{a == b} and \texttt{a is b} can give different results in Python.
\end{enumerate}


\section{Stacks}
A stack supports two operations: \texttt{push(x)} adds an element on top, and \texttt{pop()} removes and returns the top element. The function \texttt{is\_empty(S)} returns true if the stack $S$ has no elements.
\begin{lstlisting}[language=Python]
class Stack:
    def __init__(self):
        self.items = []          # the underlying list

    def push(self, x):
        self.items.append(x)

    def pop(self):
        if self.is_empty():
            raise IndexError("pop from empty stack")
        return self.items.pop()

    def is_empty(self):
        return len(self.items) == 0
\end{lstlisting}
\begin{tipbox}{Tip}
Use \texttt{collections.deque} when you need fast appends and pops from both ends. Calling \texttt{list.pop(0)} takes $O(n)$ time.
\end{tipbox}
\section{Linked lists in C}
Each node stores a value and a pointer to the next node:
\begin{lstlisting}[language=C]
struct node {
    int value;
    struct node *next;
};

int length(struct node *head) {
    int n = 0;
    for (struct node *p = head; p != NULL; p = p->next)
        n++;
    return n;
}
\end{lstlisting}
\begin{notebox}
\textbf{Note.} If \texttt{head == NULL} the list is empty and \texttt{length} returns 0. Always check for \texttt{NULL} before dereferencing a pointer.
\end{notebox}
\section{Algorithm: binary search}
\begin{tcolorbox}[colback=white,colframe=red!60!black,title={Algorithm 1: BinarySearch(A, x)}]
\begin{tabbing}
\quad\=\quad\=\quad\=\kill
1: $lo \gets 0$, $hi \gets n - 1$\\
2: \textbf{while} $lo \le hi$ \textbf{do}\\
3: \> $mid \gets \lfloor (lo + hi)/2 \rfloor$\\
4: \> \textbf{if} $A[mid] = x$ \textbf{then return} $mid$\\
5: \> \textbf{else if} $A[mid] < x$ \textbf{then} $lo \gets mid + 1$\\
6: \> \textbf{else} $hi \gets mid - 1$\\
7: \textbf{return} $-1$
\end{tabbing}
\end{tcolorbox}
Binary search runs in $O(\log n)$ time because the search interval halves at each step. The loop invariant is: if $x$ is in $A$, then $x \in A[lo..hi]$.
\begin{center}
\begin{tabular}{lll}
\textbf{Command} & \textbf{Meaning} & \textbf{Example}\\
\texttt{git status} & show changed files & \texttt{git status -s}\\
\texttt{git add} & stage changes & \texttt{git add main.py}\\
\texttt{git commit} & record a snapshot & \texttt{git commit -m "fix"}\\
\texttt{ls -la} & list all files & \texttt{ls -la /tmp}\\
\end{tabular}
\end{center}
\section{Exercises}
\begin{enumerate}
\item Write a function \texttt{reverse(head)} that reverses a linked list in place.
\item What is the output of \texttt{print([x*x for x in range(5)])}?
\item Explain why \texttt{a == b} and \texttt{a is b} can give different results in Python.
\end{enumerate}

\end{document}
